% ===========================================================================
% CHAPTER 3
% ===========================================================================

\chapter{Requirements and System Analysis}
\label{chap:requirements-system-analysis}

\section{Stakeholders and Use Cases}
\label{sec:stakeholders-use-cases}

SSL Guardian is intended primarily for certificate analysts, cybersecurity
researchers, and system operators. Analysts use the platform to examine
certificate distributions, filter records, compare global and country-level
scopes, and investigate individual certificates or shared-key groups. System
operators prepare acquisition inputs, monitor collection processes, maintain
the certificate dataset, and refresh analytical results.

The principal use cases are:

\begin{enumerate}[label=\arabic*.,leftmargin=0.42in]
  \item acquire and structure certificates presented by TLS endpoints;
  \item identify newly observed and renewed certificates through CT-assisted monitoring;
  \item analyze validity, cryptographic, issuer, SAN, geographic, and shared-key characteristics;
  \item compare Certificate Authorities using certificate-derived metrics;
  \item search, filter, and inspect certificate records through an interactive interface; and
  \item examine selected results at global and configured country-level scopes.
\end{enumerate}

\section{Functional Requirements}
\label{sec:derived-functional-requirements}

Table~\ref{tab:derived-functional-requirements} summarizes the principal
functional requirements of the developed system.

\begin{table}[htbp]
  \centering
  \small
  \setlength{\tabcolsep}{4pt}
  \caption{Functional requirements}
  \label{tab:derived-functional-requirements}
  \begin{tabularx}{\textwidth}{@{}>{\Centering\arraybackslash}p{0.10\textwidth}Y@{}}
    \toprule
    \textbf{ID} & \textbf{Requirement} \\
    \midrule
    FR-01 & The system shall accept domain datasets and maintain persistent acquisition states. \\
    FR-02 & The system shall retrieve peer certificates from TLS endpoints and convert them into structured X.509 records. \\
    FR-03 & The system shall retain domain, scan-time, scope, leaf-classification, fingerprint, parsed-certificate, and certificate-quality information. \\
    FR-04 & The system shall consume CT-derived domain observations and distinguish renewal candidates from newly observed domains. \\
    FR-05 & The system shall incorporate successfully acquired renewal and discovery results into the certificate dataset and refresh analytical results. \\
    FR-06 & The system shall compute global and country-level analytics for validity, geography, SANs, shared keys, signatures, hashes, and Certificate Authorities. \\
    FR-07 & The system shall provide searchable and filterable certificate lists, detailed records, analytical summaries, and time-based views. \\
    FR-08 & The system shall present the results through an interactive web-based dashboard. \\
    \bottomrule
  \end{tabularx}
\end{table}

\section{Non-Functional Requirements and Constraints}
\label{sec:derived-quality-attributes}

The principal non-functional requirements concern processing capacity,
recoverability, data currency, responsiveness, usability, and maintainability.
Table~\ref{tab:derived-quality-attributes} summarizes the corresponding design
mechanisms and evaluation boundaries.

\begin{table}[htbp]
  \centering
  \small
  \setlength{\tabcolsep}{4pt}
  \caption{Non-functional requirements and constraints}
  \label{tab:derived-quality-attributes}
  \begin{tabularx}{\textwidth}{@{}>{\RaggedRight\arraybackslash}p{0.21\textwidth}>{\RaggedRight\arraybackslash}p{0.40\textwidth}Y@{}}
    \toprule
    \textbf{Attribute} & \textbf{Supporting design} & \textbf{Constraint} \\
    \midrule
    Processing capacity & Concurrent acquisition, batch writes, indexes, pagination, and materialized analytics & Performance remains dependent on workload and infrastructure \\
    Recoverability & Persistent task states, retries, stale-task recovery, checkpoints, and failure logs & Stage-level recovery does not provide a system-wide transaction \\
    Data currency & Continuous CT observation and periodic renewal/discovery processing & Refresh is cycle-based rather than instantaneous \\
    Responsiveness & Indexed queries, precomputed summaries, bounded result sets, and optional caching & No fixed latency guarantee is assumed \\
    Usability and maintainability & Feature-oriented views, shared filters, typed data contracts, modular services, and reusable components & Production hardening and formal usability evaluation remain outside the project scope \\
    \bottomrule
  \end{tabularx}
\end{table}

Certificate acquisition intentionally records presented certificates even when
normal trust or hostname validation would reject the connection. The platform
is consequently an analytical measurement system rather than a certificate
validation authority. Its current operational model is most appropriate for a
controlled research environment rather than an exposed multi-user service.

\section{Technology Stack}
\label{sec:technology-stack}

Table~\ref{tab:technology-stack} presents the principal technologies and their
roles within the platform.

\begin{table}[htbp]
  \centering
  \small
  \setlength{\tabcolsep}{4pt}
  \caption{Technology stack}
  \label{tab:technology-stack}
  \begin{tabularx}{\textwidth}{@{}>{\RaggedRight\arraybackslash}p{0.22\textwidth}>{\RaggedRight\arraybackslash}p{0.31\textwidth}Y@{}}
    \toprule
    \textbf{Layer} & \textbf{Technology} & \textbf{Role} \\
    \midrule
    Web interface & Next.js, React, TypeScript, Tailwind CSS, Recharts & Interactive pages, typed client logic, responsive presentation, and charts \\
    Application service & Django & Routing, middleware, request handling, serialization, and JSON responses \\
    Data platform & MongoDB and PyMongo & Certificate storage, processing state, CT staging, and analytical results \\
    Optional cache & Redis & Scope-aware reuse of serialized analytical responses \\
    Certificate acquisition & Python networking, TLS, and concurrency libraries & Endpoint communication, task processing, and recovery control \\
    Certificate parsing & \texttt{zcertificate} and ZLint output & Structured X.509 parsing and certificate-quality findings \\
    CT processing & CertStream-compatible service and WebSocket client & CT-domain observation, renewal detection, and discovery input \\
    \bottomrule
  \end{tabularx}
\end{table}

\section{System Components and Context}
\label{sec:system-context}

The system is organized into five cooperating component groups:

\begin{enumerate}[label=\arabic*.,leftmargin=0.42in]
  \item \textbf{Certificate acquisition} retrieves, parses, enriches, and stores certificates from domain-based inputs.
  \item \textbf{CT-assisted maintenance} observes CT-derived domains, identifies renewal and discovery candidates, and updates the certificate dataset.
  \item \textbf{Analytics processing} creates indexed global and country-level summaries while retaining live access to detailed records and trends.
  \item \textbf{Application services} expose filtered certificate data and analytical results through structured request and response interfaces.
  \item \textbf{Interactive presentation} provides dashboards, charts, tables, search, filtering, and certificate drill-down.
\end{enumerate}

Internet TLS endpoints and CT logs act as external data sources. The certificate
parser provides structured X.509 and certificate-quality information, while
MongoDB forms the main persistence and integration boundary. Analysts interact
with the system through the dashboard, and operators manage acquisition and
analytical refresh activities. The experimental IP-based crawler remains a
separate acquisition capability and does not feed the primary dashboard
dataset.
